\section{Related work}
\label{sec:related}

\paragraph{Soundness of the CoqHammer translation.}
The implemented translation is described in
\cite{CzajkaKaliszyk2018}, which states plainly that it ``is
neither sound nor complete'' and confines the theoretical
discussion to \S5.6, including the guard counterexample revisited
in \cref{rem:jar-counterexample}.  The rigorous core is
\cite{Czajka2016}: a shallow embedding of proof-irrelevant Pure
Type Systems into \emph{minimal intuitionistic} first-order logic,
proved sound by a constructive, proof-theoretic reconstruction
argument over $\eta$-long normal derivations.  Our
\cref{sec:translation} complements that result with a classical,
model-theoretic soundness theorem for a CIC-style fragment with
inductive types, in which the guard predicate is interpreted as set
membership.  The two approaches are genuinely different in kind:
the reconstruction argument yields proof terms but, as
\cite[\S6]{Czajka2016} notes, is incompatible with the
model-theoretic methods of monotonicity inference, whose adaptation
to dependent type theory is posed there as an open problem.
\Cref{sec:erasure} answers that problem for the monomorphic
fragment.  The set-theoretic models we rely on are those of
Werner, Aczel and Lee--Werner
\cite{Werner1997,Aczel1998,LeeWerner2011}.

\paragraph{Type encodings and monotonicity.}
Type guards and tags for encoding sorted logics in untyped
first-order logic go back a long way; systematic soundness analyses
appear in Claessen, Lillieström and Smallbone \cite{Claessen2011},
who introduced monotonicity inference with the naked-variable
calculus and its SAT-based refinement, and in Blanchette, B\"ohme,
Popescu and Smallbone \cite{Blanchette2016}, who extended the
framework to polymorphism, introduced the lightweight and
featherweight encodings, and mechanised the monomorphic metatheory
in Isabelle/HOL \cite{BlanchettePopescu2013}.  Couchot and Lescuyer
\cite{CouchotLescuyer2007} and Leino and R\"ummer
\cite{LeinoRummer2010} studied guard- and tag-style encodings of
polymorphism for SMT and intermediate verification languages.  Our
\cref{sec:erasure} transplants the monotonicity method to the
guarded unsorted problems produced from dependent type theory; the
technical novelties forced by the setting are the sorted-core
extraction and relativisation (\cref{thm:relativisation}), the
single-domain merge in \cref{thm:erasure-main}, the explicit
inhabitation side condition (many-sorted logic assumes nonempty
sorts, CIC does not --- \cref{prop:nec-inhab}), and the internal
infinity criterion via the translation's own injectivity and
discrimination axioms (\cref{lem:ind-infinity}), which closes the
introspection soundness gap noted in
\cite[\S7.4]{Blanchette2016}.

\paragraph{Monomorphisation.}
Heuristic monomorphisation for Sledgehammer is defined in
\cite[\S5.6, \S7.1]{Blanchette2016}, with the iteration and size
bounds we adopt; its incompleteness is remarked there without
proof.  Bobot and Paskevich \cite{BobotPaskevich2011} proved that
complete monomorphisation is uncomputable for polymorphic
first-order logic and designed the symbol-discrimination
transformation \textsf{Dis}, whose keep-the-generic-formula design
mirrors our keep policy.  \Cref{thm:undecidable} adapts their
combinatory-logic reduction to $\CICm$ with a semantic notion of
consequence, adding the decidability analysis of the instance-set
fragment (Horn saturation over reflected types) in full.  Empirical
comparisons across seventeen provers \cite{Desharnais2022} found
monomorphic encodings superior even to native polymorphic prover
support (cf.\ the polymorphic extensions of Vampire
\cite{BhayatReger2020}); this motivates investing in
monomorphisation for CoqHammer rather than waiting for TF1-native
back ends.  For dependent type theory, Lean-auto
\cite{QianEtAl2025} instantiates ``constant instances'' including
dependent arguments during its preprocessing, which corresponds to
our general instantiations (\cref{rem:tier2}); its published
description argues correctness informally, while our
\cref{thm:mono-sound} gives the semantic soundness statement and
proof for the CIC fragment.  Premise selection for hammers in
dependent type theory \cite{ZhuEtAl2025} is orthogonal: it chooses
$P$, we transform it.

\paragraph{Arity and predicate optimisations.}
The flattening steps (S3)--(S4) are folklore for hammers and were
analysed for higher-order-to-first-order translations by Meng and
Paulson \cite{MengPaulson2008}; CoqHammer performs them with
bridging axioms \cite[\S5.5]{CzajkaKaliszyk2018}.  The contribution
of \cref{sec:specialization} is the uniform soundness principle
covering them together with guard specialisation and type-argument
fusion, and the faithfulness theorem for the bridge-free versions
on the first-order core.
